\chapter{Native Extensions}\label{ch:native-extensions}

\index{native extensions}
\index{extensions}
\index{.dylib}
\index{.so}

Lattice's standard library covers a broad range of tasks---files, networking, JSON, cryptography.
But eventually you will hit a wall: a C library for image processing, a database driver, or a hardware interface that does not exist in pure Lattice.
Native extensions bridge that gap.
They let you write performance-critical or platform-specific code in C, compile it into a shared library, and call it seamlessly from Lattice.

In this chapter we will learn the extension API, build an extension from scratch, survey the extensions that ship with Lattice, and understand how the runtime discovers and loads shared libraries.

\section{The Extension API}\label{sec:ext-api}
\index{extension API}
\index{lattice\_ext.h}
\index{lat\_ext\_init}
\index{lat\_ext\_register}

\subsection{The Contract}

Every native extension is a shared library (\lstinline|.dylib| on macOS, \lstinline|.so| on Linux, \lstinline|.dll| on Windows) that exports exactly one symbol: \lstinline|lat_ext_init|.
When Lattice loads your extension, it calls this function, passing an opaque \lstinline|LatExtContext|.
Your job is to call \lstinline|lat_ext_register| to register each function you want to expose:

\begin{defnbox}{Extension Initialization}
Every Lattice extension must export a function with the signature
\lstinline|void lat_ext_init(LatExtContext *ctx)|.
Inside this function, call \lstinline|lat_ext_register(ctx, name, fn)| once for each function you want to make available to Lattice code.
\end{defnbox}

The entire public API lives in a single header: \filename{include/lattice\_ext.h}.
Extensions compile against this header only---the internal \lstinline|LatValue| layout is hidden behind opaque \lstinline|LatExtValue| pointers, so your extension code never depends on Lattice internals.

\subsection{Loading an Extension from Lattice}

From Lattice code, you load an extension with \lstinline|require_ext()|:

\begin{lstlisting}[language=Lattice, caption={Loading a native extension}]
let db = require_ext("sqlite")

let conn = db["open"]("my_database.db")
let rows = db["query"](conn, "SELECT name, age FROM users WHERE age > ?", [21])

for row in rows {
  print("${row["name"]} is ${row["age"]} years old")
}

db["close"](conn)
\end{lstlisting}

The call \lstinline|require_ext("sqlite")| returns a \lstinline|Map| where each key is a function name (like \lstinline|"open"|, \lstinline|"query"|, \lstinline|"close"|) and each value is a callable native closure.
You call these functions exactly like any other Lattice function.

\begin{tipbox}{Map-Based Namespacing}
Extensions return a \lstinline|Map| rather than injecting globals.
This keeps the global namespace clean and makes it explicit which extension a function came from.
If you prefer dot-syntax, you can assign the map to a variable and use bracket notation: \lstinline|db["query"](conn, sql)|.
\end{tipbox}

\subsection{The LatExtFn Signature}

Every registered extension function has the same C signature:

\begin{verbatim}
  typedef LatExtValue *(*LatExtFn)(LatExtValue **args, size_t argc);
\end{verbatim}

The function receives an array of \lstinline|LatExtValue| pointers (one per argument) and the argument count.
It returns a heap-allocated \lstinline|LatExtValue|, or \lstinline|NULL| to return \lstinline|nil|.

Here is a minimal extension function:

\begin{verbatim}
  static LatExtValue *my_add(LatExtValue **args, size_t argc) {
      if (argc < 2) return lat_ext_error("add() expects 2 arguments");
      int64_t a = lat_ext_as_int(args[0]);
      int64_t b = lat_ext_as_int(args[1]);
      return lat_ext_int(a + b);
  }
\end{verbatim}

\section{Value Constructors and Accessors}\label{sec:ext-values}
\index{LatExtValue}
\index{extension!value constructors}
\index{extension!value accessors}

The extension API provides a complete set of constructors and accessors for working with Lattice values from C.
All values are wrapped in an opaque \lstinline|LatExtValue| pointer---you never touch raw \lstinline|LatValue| structs.

\subsection{Constructors}

\begin{table}[h]
\centering
\caption{Extension value constructors}\label{tab:ext-constructors}
\begin{tabular}{@{} l l p{5.5cm} @{}}
\toprule
\textbf{Function} & \textbf{Returns} & \textbf{Description} \\
\midrule
\lstinline|lat_ext_int(v)| & \lstinline|LatExtValue*| & Create an integer value \\
\lstinline|lat_ext_float(v)| & \lstinline|LatExtValue*| & Create a float value \\
\lstinline|lat_ext_bool(v)| & \lstinline|LatExtValue*| & Create a boolean value \\
\lstinline|lat_ext_string(s)| & \lstinline|LatExtValue*| & Create a string (copies the input) \\
\lstinline|lat_ext_nil()| & \lstinline|LatExtValue*| & Create a nil value \\
\lstinline|lat_ext_array(elems, len)| & \lstinline|LatExtValue*| & Create an array from an element array \\
\lstinline|lat_ext_map_new()| & \lstinline|LatExtValue*| & Create an empty map \\
\lstinline|lat_ext_error(msg)| & \lstinline|LatExtValue*| & Create an error value (triggers a Lattice error) \\
\bottomrule
\end{tabular}
\end{table}

All constructors return a heap-allocated \lstinline|LatExtValue*|.
The Lattice runtime deep-clones the value when receiving it, so you do not need to worry about the original being freed.

\begin{notebox}{Error Handling}
Returning \lstinline|lat_ext_error("message")| from an extension function causes Lattice to raise a runtime error.
The error message is available in \lstinline|try|/\lstinline|catch| blocks just like any other error.
\end{notebox}

\subsection{Type Queries}

Before accessing a value's contents, check its type:

\begin{verbatim}
  LatExtType lat_ext_type(const LatExtValue *v);
\end{verbatim}

The \lstinline|LatExtType| enum includes: \lstinline|LAT_EXT_INT|, \lstinline|LAT_EXT_FLOAT|, \lstinline|LAT_EXT_BOOL|, \lstinline|LAT_EXT_STRING|, \lstinline|LAT_EXT_ARRAY|, \lstinline|LAT_EXT_MAP|, \lstinline|LAT_EXT_NIL|, and \lstinline|LAT_EXT_OTHER| for types not directly representable (like closures or structs).

\subsection{Accessors}

\begin{table}[h]
\centering
\caption{Extension value accessors}\label{tab:ext-accessors}
\begin{tabular}{@{} l l p{5cm} @{}}
\toprule
\textbf{Function} & \textbf{Returns} & \textbf{Description} \\
\midrule
\lstinline|lat_ext_as_int(v)| & \lstinline|int64_t| & Extract integer value \\
\lstinline|lat_ext_as_float(v)| & \lstinline|double| & Extract float value \\
\lstinline|lat_ext_as_bool(v)| & \lstinline|bool| & Extract boolean value \\
\lstinline|lat_ext_as_string(v)| & \lstinline|const char*| & Extract string pointer (borrowed) \\
\lstinline|lat_ext_array_len(v)| & \lstinline|size_t| & Get array length \\
\lstinline|lat_ext_array_get(v, i)| & \lstinline|LatExtValue*| & Get array element (deep clone) \\
\lstinline|lat_ext_map_get(v, key)| & \lstinline|LatExtValue*| & Get map value by key \\
\lstinline|lat_ext_map_set(m, k, v)| & \lstinline|void| & Set a map entry \\
\bottomrule
\end{tabular}
\end{table}

\begin{warnbox}{Memory Management}
Values returned by \lstinline|lat_ext_array_get| and \lstinline|lat_ext_map_get| are heap-allocated deep clones.
You must call \lstinline|lat_ext_free(v)| on them when you are done.
The string returned by \lstinline|lat_ext_as_string| is borrowed---do \emph{not} free it.
\end{warnbox}

\subsection{Cleanup}

Always free extension values you create for intermediate use:

\begin{verbatim}
  void lat_ext_free(LatExtValue *v);
\end{verbatim}

The return value of your extension function is freed automatically by the runtime after it has been cloned into the VM's value space.
But any intermediate values you create (e.g., when iterating over an array argument) must be freed by your code.


\section{Building a .dylib/.so Extension}\label{sec:building-ext}
\index{extension!building}
\index{Makefile!extension}

Let us build a complete extension from scratch.
We will create a \lstinline|"math_extra"| extension that provides a fast integer exponentiation function.

\subsection{Step 1: Write the C Code}

Create a file called \lstinline|math_extra.c|:

\begin{verbatim}
  #include "lattice_ext.h"

  void lat_ext_init(LatExtContext *ctx);

  static LatExtValue *math_pow(LatExtValue **args, size_t argc) {
      if (argc < 2)
          return lat_ext_error("pow() expects 2 arguments");
      if (lat_ext_type(args[0]) != LAT_EXT_INT ||
          lat_ext_type(args[1]) != LAT_EXT_INT)
          return lat_ext_error("pow() expects integer arguments");

      int64_t base = lat_ext_as_int(args[0]);
      int64_t exp  = lat_ext_as_int(args[1]);
      int64_t result = 1;

      for (int64_t i = 0; i < exp; i++) {
          result *= base;
      }
      return lat_ext_int(result);
  }

  static LatExtValue *math_is_prime(LatExtValue **args, size_t argc) {
      if (argc < 1)
          return lat_ext_error("is_prime() expects 1 argument");
      int64_t n = lat_ext_as_int(args[0]);
      if (n < 2) return lat_ext_bool(false);
      for (int64_t i = 2; i * i <= n; i++) {
          if (n % i == 0) return lat_ext_bool(false);
      }
      return lat_ext_bool(true);
  }

  void lat_ext_init(LatExtContext *ctx) {
      lat_ext_register(ctx, "pow",      math_pow);
      lat_ext_register(ctx, "is_prime", math_is_prime);
  }
\end{verbatim}

\subsection{Step 2: Write a Makefile}

The build process is straightforward---compile the C file into a shared library, linking against the Lattice extension header:

\begin{verbatim}
  CC     = cc
  CFLAGS = -std=c11 -Wall -Wextra -fPIC -I../../include

  UNAME_S := $(shell uname -s)
  ifeq ($(UNAME_S),Darwin)
      SHARED_EXT = .dylib
      SHARED_FLAGS = -dynamiclib -undefined dynamic_lookup
  else
      SHARED_EXT = .so
      SHARED_FLAGS = -shared
  endif

  TARGET = math_extra$(SHARED_EXT)

  all: $(TARGET)

  $(TARGET): math_extra.c
  	$(CC) $(CFLAGS) $(SHARED_FLAGS) -o $@ $<

  install: $(TARGET)
  	@mkdir -p $(HOME)/.lattice/ext
  	cp $(TARGET) $(HOME)/.lattice/ext/

  clean:
  	rm -f $(TARGET)
\end{verbatim}

Key details:
\begin{itemize}
  \item \lstinline|-fPIC| generates position-independent code required for shared libraries.
  \item \lstinline|-I../../include| points to the Lattice include directory where \filename{lattice\_ext.h} lives.
  \item On macOS, \lstinline|-dynamiclib -undefined dynamic_lookup| creates a \lstinline|.dylib| that resolves Lattice symbols at load time (they are provided by the \lstinline|clat| binary).
  \item On Linux, \lstinline|-shared| creates a \lstinline|.so|.
\end{itemize}

\subsection{Step 3: Build and Install}

\begin{lstlisting}[language=Lattice, caption={Building and installing an extension}]
// From the extension directory:
// make
// make install   (copies to ~/.lattice/ext/)
\end{lstlisting}

\subsection{Step 4: Use It from Lattice}

\begin{lstlisting}[language=Lattice, caption={Using the math\_extra extension}]
let math = require_ext("math_extra")

print(math["pow"](2, 10))       // 1024
print(math["pow"](3, 5))        // 243
print(math["is_prime"](17))     // true
print(math["is_prime"](100))    // false
\end{lstlisting}

\begin{tipbox}{Returning Compound Values}
Extensions can return arrays and maps, not just scalars.
Use \lstinline|lat_ext_map_new()| and \lstinline|lat_ext_map_set()| to build structured results---database rows, image metadata, or configuration objects.
\end{tipbox}


\section{Available Extensions}\label{sec:available-extensions}
\index{extensions!available}

Lattice ships with six extensions in the \filename{extensions/} directory.
Each is a self-contained C project with its own \lstinline|Makefile|.

\subsection{ffi --- Foreign Function Interface}
\index{ffi extension}

The FFI extension lets you call functions from arbitrary shared libraries at runtime---no C code required.
It supports struct marshalling, callback trampolines, memory operations, and type signatures.

\begin{lstlisting}[language=Lattice, caption={Using the FFI extension}]
let ffi = require_ext("ffi")

// Open the C math library
let libm = ffi["open"]("/usr/lib/libm.dylib")

// Define a function signature: takes double, returns double
let sqrt_fn = ffi["sym"](libm, "sqrt", "d:d")

// Call it
let result = ffi["call"](sqrt_fn, 144.0)
print(result) // 12.0

ffi["close"](libm)
\end{lstlisting}

Core functions include \lstinline|open|, \lstinline|close|, \lstinline|sym|, \lstinline|call|, \lstinline|addr|, \lstinline|errno|, and \lstinline|strerror|.
Struct marshalling is available via \lstinline|struct_define|, \lstinline|struct_alloc|, \lstinline|struct_get|, \lstinline|struct_set|, and \lstinline|struct_to_map|.
Memory operations include \lstinline|alloc|, \lstinline|free|, \lstinline|read_i32|, \lstinline|write_i32|, and friends.

\subsection{sqlite --- SQLite Database}
\index{sqlite extension}

A complete SQLite client with parameterized queries, transactions, and metadata:

\begin{lstlisting}[language=Lattice, caption={Using the SQLite extension}]
let db = require_ext("sqlite")

let conn = db["open"]("app.db")
db["exec"](conn, "CREATE TABLE IF NOT EXISTS users (name TEXT, age INT)")
db["exec"](conn, "INSERT INTO users VALUES (?, ?)", ["Alice", 30])
db["exec"](conn, "INSERT INTO users VALUES (?, ?)", ["Bob", 25])

let rows = db["query"](conn, "SELECT * FROM users WHERE age > ?", [20])
for row in rows {
  print("${row["name"]}: ${row["age"]}")
}
// Alice: 30
// Bob: 25

db["close"](conn)
\end{lstlisting}

Functions: \lstinline|open|, \lstinline|close|, \lstinline|query|, \lstinline|exec|, \lstinline|status|, \lstinline|last_insert_rowid|.

\subsection{pg --- PostgreSQL}
\index{pg extension}

A PostgreSQL client built on \lstinline|libpq|:

\begin{lstlisting}[language=Lattice, caption={Using the PostgreSQL extension}]
let pg = require_ext("pg")

let conn = pg["connect"]("host=localhost dbname=myapp")
let rows = pg["query"](conn, "SELECT id, email FROM accounts LIMIT 10")

for row in rows {
  print("${row["id"]}: ${row["email"]}")
}

pg["close"](conn)
\end{lstlisting}

\subsection{redis --- Redis Client}
\index{redis extension}

A minimal Redis client that implements the RESP protocol over raw TCP sockets---no external dependencies:

\begin{lstlisting}[language=Lattice, caption={Using the Redis extension}]
let redis = require_ext("redis")

let conn = redis["connect"]("127.0.0.1", 6379)
redis["set"](conn, "greeting", "Hello from Lattice!")
let val = redis["get"](conn, "greeting")
print(val) // Hello from Lattice!

redis["incr"](conn, "counter")
redis["incr"](conn, "counter")
print(redis["get"](conn, "counter")) // 2

redis["close"](conn)
\end{lstlisting}

Functions: \lstinline|connect|, \lstinline|close|, \lstinline|command|, \lstinline|get|, \lstinline|set|, \lstinline|del|, \lstinline|exists|, \lstinline|expire|, \lstinline|keys|, \lstinline|incr|, \lstinline|lpush|, \lstinline|lrange|, \lstinline|publish|, \lstinline|ping|.

\subsection{image --- Image Metadata and Operations}
\index{image extension}

Reads image headers directly (PNG, JPEG, GIF, BMP, WebP) without external library dependencies.
Uses macOS \lstinline|sips| for resize, convert, and thumbnail operations:

\begin{lstlisting}[language=Lattice, caption={Using the image extension}]
let img = require_ext("image")

let info = img["info"]("photo.jpg")
print("Format: ${info["format"]}")   // jpeg
print("Size: ${info["width"]}x${info["height"]}")

// Resize (macOS only, uses sips):
img["resize"]("photo.jpg", 800, 600, "photo_small.jpg")
\end{lstlisting}

\subsection{websocket --- WebSocket Client and Server}
\index{websocket extension}

Implements RFC 6455 WebSocket framing with support for both client and server roles, ping/pong, fragmented messages, and timeouts:

\begin{lstlisting}[language=Lattice, caption={Using the WebSocket extension}]
let ws = require_ext("websocket")

// Client connection
let conn = ws["connect"]("ws://echo.websocket.org")
ws["send"](conn, "Hello, WebSocket!")
let response = ws["recv"](conn)
print(response) // Hello, WebSocket!
ws["close"](conn)
\end{lstlisting}

Functions include \lstinline|connect|, \lstinline|connect_auto| (auto-reconnect), \lstinline|listen|, \lstinline|accept|, \lstinline|send|, \lstinline|recv|, \lstinline|recv_timeout|, \lstinline|send_binary|, \lstinline|close|, \lstinline|status|, \lstinline|ping|, and \lstinline|set_timeout|.


\section{Extension Search Paths}\label{sec:ext-search-paths}
\index{extension!search paths}
\index{LATTICE\_EXT\_PATH}

When you call \lstinline|require_ext("name")|, the runtime searches for the shared library in the following order:

\begin{enumerate}
  \item \lstinline|./extensions/name.dylib| (or \lstinline|.so| or \lstinline|.dll|) --- the local project's extensions directory.
  \item \lstinline|./extensions/name/name.dylib| --- a subdirectory within the local extensions directory (the default layout for extensions that ship with the repository).
  \item \lstinline|~/.lattice/ext/name.dylib| --- the user's global extension directory.
  \item \lstinline|$LATTICE_EXT_PATH/name.dylib| --- a custom path set via the \lstinline|LATTICE_EXT_PATH| environment variable.
\end{enumerate}

The first match wins. If no library is found, \lstinline|require_ext| raises an error listing all the paths that were searched.

\begin{lstlisting}[language=Lattice, caption={Extension search path behavior}]
// If you have an extension in your project:
//   ./extensions/sqlite/sqlite.dylib   <-- found at path #2

// If you installed globally:
//   ~/.lattice/ext/sqlite.dylib        <-- found at path #3

// If you set a custom path:
//   LATTICE_EXT_PATH=/opt/lattice/ext clat my_program.lat
//   /opt/lattice/ext/sqlite.dylib      <-- found at path #4
\end{lstlisting}

\begin{tipbox}{Global Installation}
Each extension's \lstinline|Makefile| includes a \lstinline|make install| target that copies the built library to \lstinline|~/.lattice/ext/|. This makes the extension available to all your Lattice projects without copying files:
\begin{verbatim}
  cd extensions/sqlite && make && make install
\end{verbatim}
\end{tipbox}

\subsection{Security Considerations}
\index{extension!security}

Native extensions execute arbitrary C code with full process permissions.
The runtime validates that the extension name does not contain path separators (\lstinline|/| or \lstinline|\|) to prevent path traversal attacks, but it does not sandbox the loaded code.

\begin{warnbox}{Trust Your Extensions}
Only load extensions from sources you trust.
A malicious extension has the same access as the Lattice process itself---it can read files, make network connections, or modify memory.
\end{warnbox}

\subsection{WASM Limitation}

When Lattice is compiled to WebAssembly (via Emscripten), native extensions are not available.
The \lstinline|require_ext| function returns an error explaining that dynamic loading is unsupported in the WASM environment.


\section*{Exercises}

\begin{enumerate}
  \item Build the \lstinline|math_extra| extension described in \cref{sec:building-ext}. Load it from a Lattice program and verify that \lstinline|pow(2, 20)| returns \lstinline|1048576|.

  \item Add a \lstinline|gcd| function to the \lstinline|math_extra| extension that computes the greatest common divisor of two integers using Euclid's algorithm. Register it in \lstinline|lat_ext_init| and test it from Lattice.

  \item Write an extension function that accepts a Lattice array of strings and returns them sorted alphabetically using C's \lstinline|qsort|. Use \lstinline|lat_ext_array_len|, \lstinline|lat_ext_array_get|, and \lstinline|lat_ext_as_string| to extract the strings, sort them, and return a new array via \lstinline|lat_ext_array|.

  \item The extension API uses opaque \lstinline|LatExtValue| pointers to hide the internal \lstinline|LatValue| layout. What are the advantages of this design over exposing the raw struct? What are the disadvantages (hint: think about performance)?

  \item (Challenge) Write an extension that wraps a third-party C library of your choice (e.g., \lstinline|zlib| for compression, \lstinline|curl| for HTTP requests). Include proper error handling for all failure cases.
\end{enumerate}

\bigskip
\begin{center}\rule{0.5\linewidth}{0.4pt}\end{center}

\textbf{What's Next}\quad
We have seen how to extend Lattice with C code. But Lattice also provides powerful tools for inspecting and manipulating programs from \emph{within} the language itself.
In \cref{ch:metaprogramming}, we explore runtime code introspection, the \lstinline|tokenize()| function, struct reflection, and the dynamic features that make Lattice a surprisingly flexible metaprogramming platform.
